\cotefile{60}
\scanpage{2}
Notes sur \uline{l'équation \unsure{fonctionnelle|fonctionelle}}
% archive sheet number: 1

\scanpage{3}
Si $f$ est un endomorphisme d'un vectoriel $E$ de dim finie $n$ sur le corps $k$, on pose
\[
\begin{aligned}
P_f(t)=\det(1-ft)&=\prod_i(1-\lambda_i t)\\
&=\sum_i(-1)^i\bigl(\operatorname{Tr}\lambda^i(f)\bigr)t^i
\end{aligned}
\]
les $\lambda_i$ étant les val propres dans une \unsure{extension|clôture} alg. close de $k$.

Donc les racines de $P_f(t)$ sont les inverses des valeurs propres non nulles de $f$.

$P_f(t)$ pour $(E,f)$ variable est ``additif'' en $E$ (\unsure{suite exacte|suites exactes} en $E$ \unsure{transformée|transformé} en produit des polynômes). Il en est donc de $=$ de
\[
\begin{aligned}
Z_f(t)=\frac{1}{P_f(t)}
 &=\prod_i\frac{1}{1-\lambda_i t}
 =\prod_i\sum_{\unsure{\nu|\kappa}}\lambda_i^\nu t^\nu\\
 &=\sum_\nu t^\nu\sum_{\nu_1+\cdots+\nu_n=\nu}\lambda_1^{\unsure{\nu_1|\nu_n}}\cdots\lambda_n^{\nu_n}
 =\sum_\nu\operatorname{Tr}\bigl(\operatorname{Sym}^\nu(f)\bigr)t^\nu
\end{aligned}
\]

D'autre part, $P_f(t)$ transforme le produit tensoriel en $(E,f)$ en le ``produit'' $\ast$ des séries formelles commençant par 1, où on a posé
\[
(1-at)\ast(1-bt)=1-abt
\]
[et distributivité de $\ast$ p.r. à\ $\cdot$, etc\ldots]
\margin{\diagram{\deleted{$\begin{array}{ccc}
p-1 & \underset{\unsure{\leftrightarrow|{-}}}{p} & p+1\\
\underset{\rightarrow}{-2} & -1 & \underset{\unsure{\leftarrow|\leftrightarrow}}{0}\\
& p,\ p+1 & \unsure{p+2|\rho+2}\\
p-1 & 0 & p+1
\end{array}$}}}
% archive sheet number: 2

\scanpage{4}
\mspage{- 57 -}
n\textsuperscript{o}357
\[
f(x)\mathrel{\handwritten{\leq}}\handwritten{\sum}_{i=1}^{n}a_i\bigl[\varphi(s_i)+\handwritten{\eta}\bigr]g(s_i^{-1}x)
\]
et de même pour $f'$, et par suite :
\[
(f:g)+(f':g)\mathrel{\handwritten{\leq}}\handwritten{\sum}_{i=1}^{n}a_i\bigl[\varphi(s_i)+\varphi'(s_i)+2\handwritten{\eta}\bigr]\mathrel{\handwritten{\leq}}(1+2\handwritten{\eta})\handwritten{\sum}_{i=1}^{n}a_i
\]
puisque $\varphi+\varphi'\mathrel{\handwritten{\leq}}1$. En appliquant la définition de $[F:g]$, puis (ii) et (vii), on en conclut :
\[
(f:g)+(f':g)\mathrel{\handwritten{\leq}}(1+2\handwritten{\eta})(F:g)\mathrel{\handwritten{\leq}}(1+2\handwritten{\eta})\bigl[(f+f':g)+\handwritten{\delta}(h:g)\bigr]\ ;
\]
appliquant (v) à $f+f'$, $h$ et $g$, cela donne :
\[
(f:g)+(f':g)\mathrel{\handwritten{\leq}}(f+f':g)+2\handwritten{\eta}(f+f':h).(h:g)+(1+2\handwritten{\eta})\handwritten{\delta}.(h:g)\ .
\]
Si donc $f$, $f'$ et $h$ sont donnés comme dans (viii), il suffit de choisir $\handwritten{\delta}$ et $\handwritten{\eta}$ assez petits, puis de choisir $V$ comme il a été dit, pour qu'on ait l'inégalité formulée dans (viii).

Soit maintenant $\handwritten{\unsure{\mathfrak{F}|\mathscr{F}}}$ un ultrafiltre sur $\handwritten{\unsure{\mathcal{K}|\mathscr{K}}}^{\ast}_{+}$, plus fin que le filtre engendré par la réunion des $\handwritten{\mathcal{K}}^{\ast}_{+}(V)$ quand $V$ parcourt un système fondamental de voisinages compacts de $e$. D'autre part, fixons $f_{\unsure{o|0}}\mathrel{\handwritten{\in}}\handwritten{\mathcal{K}}^{\ast}_{+}$ ; posons, pour $f\mathrel{\handwritten{\in}}\handwritten{\mathcal{K}}^{\ast}_{+}$, $g\mathrel{\handwritten{\in}}\handwritten{\mathcal{K}}^{\ast}_{+}$ :
\[
I_g(f)=\frac{(f:g)}{(f_o:g)}\ ,
\]
puis $I(f)=\lim_{\handwritten{\mathfrak{F}}}I_g(f)$. D'après (vi), $I(f)$ est dans l'intervalle $[1/(f_o:f),(f:f_o)]$. D'après (ii), on a $I(cf)=c.I(f)$ pour $c>0$. D'après (vii), on a $I(f+f')\mathrel{\handwritten{\leq}}I(f)+I(f')$. D'après (viii), on a $I(f+f')\mathrel{\handwritten{\geq}}I(f)+I(f')-\varepsilon I(h)$ quel que soit $\varepsilon\mathrel{\handwritten{\geq}}0$ si $h$ est $\mathrel{\handwritten{\geq}}1$ sur le support de $f+f'$ ; comme on peut toujours choisir $h$ ainsi dans $\handwritten{\mathcal{K}}^{\ast}_{+}$, il s'ensuit que $I(f+f')=I(f)+I(f')$. Alors, d'après le Chap. II, \handwritten{\S}\ 2, n\textsuperscript{o}1, prop. 2, $I$ se prolonge (d'une manière et d'une seule) à une forme linéaire positive sur $\handwritten{\mathcal{K}}$, c'est-à-dire qu'il y a une mesure positive $\alpha$ sur $G$ telle que $I(f)=\handwritten{\int} f\,d\alpha$ pour $f\mathrel{\handwritten{\in}}\handwritten{\mathcal{K}}^{\ast}_{+}$. D'après (i), on a $I(\handwritten{\gamma}(\varepsilon_s)f)=I(f)$ quels que
% archive sheet number: 3

\scanpage{5}
\[
\left\{\begin{aligned}
P_{f\otimes\unsure{f'|f}}(t)&=P_f(t)\ast P_{f'}(t)\\
L_{f\otimes f'}(t)&=\frac{1}{L_f(t)\ast L_{f'}(t)}\ ,\qquad
L_{f_1\otimes\cdots\otimes f_\nu}(t)=\left(\prod L_{f_i}(t)\right)^{(-1)^{\nu+1}}
\end{aligned}\right.
\]

De $=$, pour ``la'' $\lambda$-structure de $1+k[[t]]\unsure{+|{}^{+}}$ où $\deleted{\unsure{\psi|\varphi|\lambda^p}}\inserted{\lambda^p}$ transforme $\prod_{1}^{n}(1-\lambda_i t)$ en $\prod_{i_1<\cdots<i_p}(1-\lambda_{i_1}\cdots\lambda_{i_p}t)$, on aura
\[
\left\{\begin{aligned}
P_{\lambda^i f}(t)&=\lambda^i P_f(t)\ ,\qquad\text{donc}\\
\frac{1}{L_{\lambda^i f}(t)}&=\lambda^i\frac{1}{L_{\deleted{\unsure{\lambda^i|\lambda}}f}(t)}
\end{aligned}\right.
\]
i.e.
\[
\begin{aligned}
-\sum\deleted{\unsure{1|-1|(-1)^i}}\inserted{(-1)^i}\bigl[L_{\lambda^i f}(t)\bigr]s^i
 &=\lambda^\ast\bigl([-L_f(t)]\bigr)(-s)\\
 &=\frac{1}{\lambda^\ast[L_f(t)](-s)}\deleted{\unsure{(-s)|(s)}}\\
\deleted{\Bigl(\sum[\lambda^i(L_f(t))]s^i\Bigr)\Bigl(\sum[L_{\lambda^i(f)}(t)]\unsure{s^i|s^1}\Bigr)=-1=\Bigl[\frac{1}{1-t}\Bigr]}&\\
 &=\sum\sigma^i\bigl(L_f(t)\bigr)s^i
\end{aligned}
\]
\[
\left\{\begin{aligned}
L_{\lambda^i(f)}(t)&=\bigl(\sigma^i L_f(t)\bigr)^{(-1)^{i+1}}\\
L_{\sigma^i(f)}(t)&=\bigl(\lambda^i\unsure{(L_f(t)|L_f(t)}\bigr)^{(-1)^{i+1}}
\end{aligned}\right.
\]
% archive sheet number: 4

\scanpage{6}
\mspage{- 56 -}
n\textsuperscript{o}357
\[
0<\frac{1}{(f_{\unsure{o|0}}:f)}\mathrel{\handwritten{\leq}}\frac{(f:g)}{(f_o:g)}\mathrel{\handwritten{\leq}}(f:f_o)\ .
\]
\begingroup\csname tagsleft@true\endcsname\[
(f+f':g)\mathrel{\handwritten{\leq}}(f:g)+(f':g).\tag*{(vii)}
\]\endgroup

(viii) \quad si $f$, $f'$, $h$ sont dans $\handwritten{\unsure{\mathcal{K}|\mathscr{K}}}_+$, si $h\mathrel{\handwritten{\geq}}1$ sur le support de $f+f'$, et si $\varepsilon>0$, il y a un voisinage compact $V$ de $e$ tel que l'on ait, pour tout $g\mathrel{\handwritten{\in}}\handwritten{\mathcal{K}}^{\ast}_{+}(V)$ :
\[
(f+f':g)\mathrel{\handwritten{\geq}}(f:g)+(f':g)-\varepsilon(h:g)\ .
\]

Les trois premières sont évidentes ; (iv) aussi, puisque $\sup(\handwritten{\gamma}(\mu)g)\mathrel{\handwritten{\leq}}\lVert\mu\rVert\sup(g)$ ; (v) résulte aussitôt de la formule $\lVert\mu\ast\mu'\rVert={ }=\lVert\mu\rVert.\lVert\mu'\rVert$ (par la vertu du Chap. VII, \handwritten{\S}\ 1, n\textsuperscript{o}1, ou, moins savamment et plus simplement dans le cas qui nous occupe ici, par un calcul direct trivial) et de ce qu'on a, si $f\mathrel{\handwritten{\leq}}\handwritten{\gamma}(\mu)g$ et $g\mathrel{\handwritten{\leq}}\handwritten{\gamma}(\mu')h$ :
\[
f\mathrel{\handwritten{\leq}}\handwritten{\gamma}(\mu)\handwritten{\gamma}(\mu')h=\handwritten{\gamma}(\mu\ast\mu')h\ .
\]

Si on applique (v) à $f_o$, $f$, $g$ d'une part, et à $f$, $f_o$, $g$ d'autre part, on obtient (vi) ; (vii) est évident. Enfin, soient $f$, $f'$, $h$ comme il est dit dans (viii). Posons $F=f+f'+\handwritten{\delta} h$, avec $\handwritten{\delta}>0$ ; les fonctions $\varphi$, $\varphi'$ qui coïncident respectivement avec $f/F$ et $f'/F$ sur le support de $f+f'$, et qui sont $0$ en dehors de celui-ci, sont dans $\handwritten{\mathcal{K}}$, donc uniformément continues ; pour tout $\handwritten{\eta}>0$, on peut donc choisir un voisinage compact $V$ de $e$ tel que $|\varphi(x)-\varphi(s)|$ et $|\varphi'(x)-\varphi'(s)|$ soient $\mathrel{\handwritten{\leq}}\handwritten{\eta}$ pour $s^{-1}x\mathrel{\handwritten{\in}}V$ ; alors, quelle que soit $g\mathrel{\handwritten{\in}}\handwritten{\mathcal{K}}^{\ast}_{+}(V)$, on aura
\[
\varphi(x)g(s^{-1}x)\mathrel{\handwritten{\leq}}\unsure{[|(}\varphi(s)+\handwritten{\eta}]g(s^{-1}x)
\]
quels que soient $x$ et $s$, et de même pour $\varphi'$. Supposons maintenant qu'on ait, pour tout $x$ :
\[
F(x)\mathrel{\handwritten{\leq}}\handwritten{\sum}_{i=1}^{n}a_i g(s_i^{-1}x)\ .
\]
Puisque $f=\varphi F$, on en conclut :
% archive sheet number: 5

\scanpage{7}
Considérons le cas d'un \uline{automorphisme} $f$, et voyons quel est le passage au \unsure{contragrédient|contragredient}
\[
\begin{aligned}
P_{\check f}(t)&=\det(1-\check f t)=\prod(1-\lambda_i^{-1}\deleted{\unsure{f}}t)\\
&\deleted{=\prod\lambda_i^{-1}(\unsure{\lambda_i-t|\bar\lambda_i-t})}\\
&=\prod_i(-\lambda_i^{-1}t)(1-\lambda_i t^{-1})\\
&=\frac{(-1)^nt^n}{\prod\lambda_i}\deleted{\unsure{\det}}\prod(1-\lambda_i t^{-1})\\
&=\frac{(-1)^nt^n}{\det f}P_f(\unsure{t^{-1}|-t^{-1}})\qquad\text{donc}
\end{aligned}
\]
\[
\deleted{P_{\check f}(t)=\deleted{\frac{(-1)^nt^n}{\det f}}\unsure{\text{``}\varepsilon(f)\text{''}|{\varepsilon(f)}}t^n P_f(t^{-1})\qquad\varepsilon(f)=}
\]
\[
\boxed{P_{\check f}(t)=(-t)^{\varepsilon(f)}\det f^{-1}\ P_f(t^{-1})}\qquad\text{i.e}
\]
\[
\boxed{L_{\check f}(t)=(-t)^{-\varepsilon(f)}\det f\ L_f(t^{-1})}
\]
Bien entendu, le deuxième \unsure{mb} n'a pas de sens quand on remplace $P_f^{\unsure{=|-}}(t)$ par une série formelle quelconque : il est essentiel de se borner aux séries formelles qui sont \unsure{produits|justement} alternés finis de \uline{polynômes} [à terme constant 1, à terme dominant inversible -- cette dernière condition n'étant utile que si $k$ est un anneau général \ldots].
% archive sheet number: 6

\scanpage{8}
\mspage{- 55 -}
n\textsuperscript{o}357

[N.B. -- Ce n'est pas par amour des méthodes éculées et pisseuses que le rédacteur, dans la construction de la mesure de Haar, est revenu aux mesures à support \uline{fini} au lieu des mesures à support \uline{compact} des rédactions précédentes ; c'est en vertu du rasoir d'Occam (``entitates non sunt multiplicandae''), et pour pouvoir laisser reposer en paix dans la Remarque de la fin du Chap. VII, \handwritten{\S}\ 2, n\textsuperscript{o}2, la théorie générale de l'opération $\mu\ast f$, ou $\handwritten{\gamma}(\mu)f$, lorsqu'on ne s'est pas donné de jauge ; ce n'est évidemment pas difficile (cf. loc. cit), mais ça paraît bien superflu \ldots]

On écrira $\handwritten{\unsure{\mathcal{K}|\mathscr{K}}}$, $\handwritten{\mathcal{K}}(K)$, etc. au lieu de $\handwritten{\mathcal{K}}(G)$, $\handwritten{\mathcal{K}}(G,K)$, etc. ; on notera $\handwritten{\mathcal{K}}_+$ l'ensemble des fonctions $f\mathrel{\handwritten{\in}}\handwritten{\mathcal{K}}$ à valeurs $\mathrel{\handwritten{\geq}}0$, $\handwritten{\mathcal{K}}^{\ast}_{+}$ l'ensemble des fonctions $g\mathrel{\handwritten{\in}}\handwritten{\mathcal{K}}_+$ autres que $0$, et $\handwritten{\mathcal{K}}^{\ast}_{+}(K)$, si $K$ est une partie compacte de $G$, l'ensemble des fonctions $g\mathrel{\handwritten{\in}}\handwritten{\mathcal{K}}^{\ast}_{+}$ à support contenu dans $K$. Pour $f\mathrel{\handwritten{\in}}\handwritten{\mathcal{K}}_+$, $g\mathrel{\handwritten{\in}}\handwritten{\mathcal{K}}^{\ast}_{+}$, désignons par $A(f,g)$ l'ensemble des mesures positives $\mu\mathrel{\handwritten{\ne}}0$, à support fini, telles que $f\mathrel{\handwritten{\leq}}\handwritten{\gamma}(\mu)g$. \uline{Cet ensemble n'est pas vide}. En effet, soient $K$ le support de $f$, et $U$ un ouvert non vide tel que $m=\inf_U(g)>0$ ; soit $M=\sup_K(f)$. On peut recouvrir $K$ par un nombre fini d'ensembles $s_i U$ ; alors la mesure $(M/m)\handwritten{\sum}\varepsilon_{s_i}$ est dans $A(f,g)$. Posons maintenant :
\[
(f:g)=\inf_{\mu\mathrel{\handwritten{\in}} A(f,g)}\lVert\mu\rVert\ .
\]

Il est clair qu'on a $0\mathrel{\handwritten{\leq}}\quad(f:g)<+\handwritten{\infty}$. De plus, ce symbole a les propriétés suivantes :

(i) \quad $(\handwritten{\gamma}(\varepsilon_s)f:g)=(f:g)$ quel que soit $s\mathrel{\handwritten{\in}}G$.

(ii) \quad $(cf:g)=c.(f:g)$ quel que soit $c\mathrel{\handwritten{\geq}}0$.

(iii) \quad $f\mathrel{\handwritten{\leq}} f'$ entraîne $(f:g)\mathrel{\handwritten{\leq}}(f':g)$.

(iv) \quad $(f:g)\mathrel{\handwritten{\geq}}(\sup f)/(\sup g)$.

(v) \quad $(f:h)\mathrel{\handwritten{\leq}}(f:g).(g:h)$ pour $g\mathrel{\handwritten{\in}}\handwritten{\mathcal{K}}^{\ast}_{+}$, $h\mathrel{\handwritten{\in}}\handwritten{\mathcal{K}}^{\ast}_{+}$.

(vi) \quad si $f$, $f_{\unsure{o|0}}$, $g$ sont dans $\handwritten{\mathcal{K}}^{\ast}_{+}$, on a :
% archive sheet number: 7

\scanpage{9}
Si $E$ est un \deleted{motif}\inserted{$\unsure{Q|\mathbb{Q}}_{\ell}$-faisceau} sur $\unsure{\operatorname{Spec}|\operatorname{Sp}}\unsure{(\mathbb{F}_p,)|(\mathbb{F}_p),}$ \deleted{\unsure{sur|dans}} on trouve donc
\[
L_{\check E}(t)=(-t)^{-n}\det f_E\ L_E(t^{-1})
\]
où
\[
\left\{\begin{aligned}
n&=\operatorname{rang} E
  &&\text{et s'il s'agit d'un \uline{motif} \inserted{de poids $\rho$}}\\
 &&&\text{satisfaisant aux conj. de Weil}\\
\det f_E&=\deleted{\unsure{\varepsilon}}\inserted{\varepsilon(E)}\deleted{\unsure{\pm|+}}p^{\frac{n\rho}{2}}
\end{aligned}\right.
\]
avec $\varepsilon=\pm1$ [dépendant du motif $\bigwedge^{\unsure{n|{-}}}E$, \unsure{savoir} si $f$ y opère \unsure{positivement|parfaitement} ou non] ; noter que $\frac{n\rho}{2}$ est tjrs entier (i.e. $\rho$ impair $\Longrightarrow$ $n$ pair).

On trouve donc
\[
L_{\check E}(t)=\varepsilon(E)(-p^{\frac{\rho}{2}}/t)^n L_E(t^{-1})\ .
\]

D'autre part, par Lefschetz, on aura
\[
\check E\simeq E(\rho)
\]
d'où
\[
\begin{aligned}
L_{\check E}(t)&=\bigl(L_E(t)\ast L_{\mathbb{Q}_{\ell}(\rho)}(t)\bigr)^{-1}\\
&=\left(L_E(t)\ast\frac{1}{1-p^{-\rho}t}\right)^{-1}=L_E(t)\ast(1-p^{-\rho}t)
\end{aligned}
\]
\[
L_{\check E}(t)=L_E\left(\frac{t}{\unsure{p\rho|p^\rho}}\right)
\]
% archive sheet number: 8

\scanpage{10}
CHAPITRE VIII

CONSTRUCTION de M\unsure{E|\={E}}SURES INVARIANTES

\handwritten{\S} 1. \uline{Construction d'une mesure de Haar}.

\uline{Dans ce} \handwritten{\S}, $G$ \uline{désignera un groupe localement compact}.

1. \uline{Le théorème d'existence et d'unicité}.

Rappelons (Chap. VII, \handwritten{\S} 2, n\textsuperscript{o}4) qu'on appelle \uline{mesure de Haar} sur $G$ toute mesure positive $\alpha\mathrel{\handwritten{\ne}}0$ sur $G$, invariante à gauche, et qu'une telle mesure est alors une jauge propre à gauche et à droite, satisfaisant à une formule de la forme
\[
d\alpha(axb)=\handwritten{\Delta}_G(b)d\alpha(x)
\]
quels que soient $a\mathrel{\handwritten{\in}}G$, $b\mathrel{\handwritten{\in}}G$, $\handwritten{\Delta}_G$ étant un caractère de $G$ ; on a alors
\[
d\alpha(x^{-1})=\handwritten{\Delta}_G(x^{-1})d\alpha(x)\ .
\]
Si $\handwritten{\Delta}_G=1$, $\alpha$ est invariante à gauche et à droite ; on dit que $\alpha$ est \uline{bi-invariante} ; il en est évidemment ainsi si $G$ est commutatif.

THEOREME 1.- \uline{Sur tout groupe localement compact} $G$, \uline{il existe une mesure de Haar}, \uline{et}, \uline{à un facteur constant près}, \uline{il n'en existe qu'une}.

L'unicité a déjà été démontrée (Chap. VII, \handwritten{\S} 2, n\textsuperscript{o}4, th.1) ; on va démontrer l'existence par une construction explicite. Introduisons d'abord une notation. Soit $\mu=\handwritten{\sum}_{i=1}^{n}a_i\varepsilon_{s_i}$ une mesure de support fini $\handwritten{\{}s_1,\ldots,s_n\handwritten{\}}$ sur $G$ ; alors, pour $f\mathrel{\handwritten{\in}}\handwritten{\unsure{\mathcal{K}|\mathscr{K}}}(G)$, on désignera par $\handwritten{\gamma}(\mu)f$ la fonction $f'\mathrel{\handwritten{\in}}\handwritten{\mathcal{K}}(G)$ définie par
\[
f'(x)=\handwritten{\sum}_{i=1}^{n}a_i f(s_i^{-1}x)\ ;
\]
on notera que, dès qu'on aura construit une mesure de Haar, on aura le droit, en vertu des définitions du Chap. VII, \handwritten{\S} 2, d'écrire $\mu\ast f$ au lieu de $\handwritten{\gamma}(\mu)f$ (cf. la Remarque à la fin du Chap. VII, \handwritten{\S} 2, n\textsuperscript{o}2). L'application $(\mu,f)\mathrel{\handwritten{\longrightarrow}}\handwritten{\gamma}(\mu)f$ ainsi définie est évidemment bilinéaire ; comme on a $\handwritten{\gamma}(\varepsilon_{st})=\handwritten{\gamma}(\varepsilon_s)\handwritten{\gamma}(\varepsilon_t)$, on en conclut qu'on a $\handwritten{\gamma}(\mu\ast\handwritten{\unsure{\nu|v}})=\handwritten{\gamma}(\mu)\handwritten{\gamma}(\handwritten{\nu})$. De plus, $\handwritten{\gamma}(\mu)f$ est $\mathrel{\handwritten{\geq}}0$ si $\mu$ et $f$ sont $\mathrel{\handwritten{\geq}}0$ ; et, si $\mu$ et $f$ sont $\mathrel{\handwritten{\geq}}0$, on ne peut avoir $\handwritten{\gamma}(\mu)f=0$ que si $\mu=0$ et $f=0$.
% archive sheet number: 9

\scanpage{11}
On trouve donc
\[
\boxed{L_E\left(\deleted{\unsure{t|t^{-1}}}\frac{1}{t\unsure{p\rho|p^\rho}}\right)
=\deleted{\unsure{(-t)^n|(-1/t)^n}}\det f_E\,(-t)^n L_E(t)}
\]
[$\rho$ étant le poids de $E$, qui s'introduit via
\[
\check E\simeq E(\rho)\ \bigr]
\]

\noindent\rule{\linewidth}{0.35pt}\par
\centerline{\rule{0.58\linewidth}{0.35pt}}

Soit $E$ un $\unsure{\mathbb{Q}|Q}_{\ell}$-faisceau constructible sur $X$ \deleted{\unsure{compactifiable}}\inserted{propre} sur $\mathbb{F}_p$, alors on aura
\[
\deleted{H_{\unsure{c|{}}}^i(X,DE)\unsure{,|{}}\simeq D(H^{-i}(X,E))}
\]
\margin{par le th. de dualité globale}
\deleted{d'où dans}
\[
\deleted{\chi(DE)=\sum(-1)^iH^i(X,D(E))\simeq D(\chi(E))}
\]
\[
\mathbb{R}f_\ast(DE)=D(\mathbb{R}f_\ast(E))
\]
d'où
\begingroup\csname tagsleft@true\endcsname\[
\left\{\begin{aligned}
&\boxed{L_{DE}(t)=(-t)^{-\chi(E)}\delta(E)\,L_E(t^{-1})}\\
&\chi(E)=\operatorname{rang}\mathbb{R} f_\ast E=\sum(-1)^i\operatorname{rg} H^i(\bar X,\bar E)\\
&\delta(E)=\prod_i\bigl(\text{dét}\ f_{H^i(\bar X,\bar F)}\bigr)^{(-1)^i}=
\varepsilon(E)p^{\sum_i(-1)^i\deleted{\unsure{\rho/2}}\frac{i+\rho}{2}b_i}
\end{aligned}\right.\tag*{\raisebox{1.35\baselineskip}{(1)}}
\]\endgroup
\margin{invariant géométrique}
\margin{invariant arithmétique}
\margin{avec les conjectures \unsure{standart|standard} si $X$ est \deleted{\unsure{projectif|projective}}\inserted{lisse} \unsure{purement de dim $n$} \deleted{\unsure{lisse,}} et $F$ $c^{t}$. \unsure{tordu|tordue} purement de poids $\rho$}
% slanted
\margin{$b_i=\operatorname{rang} H^i(\bar X,\bar F)$}
% archive sheet number: 10

\scanpage{12}
\mspage{- 2 -}
n\textsuperscript{o} 398

PROPOSITION 2.- \uline{Si} $E'\longrightarrow E\longrightarrow E''$ \uline{est une suite exacte de} $A$-\uline{modules de longueur finie}, \uline{la suite} $E^{\prime\prime\ast}\longrightarrow E^\ast\longrightarrow E^{\prime\ast}$ \uline{est exacte}.

Soit d'abord $0\longrightarrow L\longrightarrow M\longrightarrow N\longrightarrow0$ une suite exacte de modules de longueur finie et montrons que $0\longrightarrow N^\ast\longrightarrow M^\ast\longrightarrow L^\ast$ est exacte. C'est en tout cas vrai pour $0\longrightarrow N^\ast\longrightarrow M^\ast\longrightarrow L^\ast$, d'après les propriétés générales de Hom. Comptant les longueurs, on en déduit que $M^\ast\longrightarrow L^\ast$ est surjectif, d'où le résultat.

Le cas d'une suite exacte qqe $E'\longrightarrow E\longrightarrow E''$ se ramène au précédent en dévissant à la façon habituelle, (Chap. II, 3e éd., p. 21). Soit $E$ un $A$-module. Si $F\subset E$, on note $F^{\unsure{o|\circ|0}}$ son orthogonal dans $E^\ast$, i.e. etc \ldots

PROPOSITION 3.- \uline{Soit} $E$ \uline{un} $A$-\uline{module de longueur finie}. \uline{La correspondance qui à} $F\subset E$ \uline{fait correspondre} $F^o$ \uline{est une bijection de l'ensemble des sous-modules de} $E$ \uline{sur celui des sous-modules de} $E^\ast$ ; \uline{cette bijection transforme somme en intersection et intersection en somme}.

On utilise la suite exacte $0\longrightarrow(E/F)^\ast\longrightarrow E^\ast\longrightarrow F^\ast\longrightarrow0$, qui montre que $F^\ast$ s'identifie à $E^\ast/F^o$ et $(E/F)^\ast$ à $F^o$. On en déduit que l'injection $F\longrightarrow F^{oo}$ s'identifie à l'injection $F\longrightarrow F^{\ast\ast}$, qui est un isomorphisme. D'où, etc \ldots

\uline{Exemple}.- $A={}$\uuline{$Z$}, $K/A={}$\uuline{$Q$}$/$\uuline{$Z$} ; on a une dualité pour les groupes abéliens finis.

\centerline{\texttt{-:-:-:-}}
% archive sheet number: 11

\scanpage{13}
% archive sheet number: 12

Dans ce cas favorable
\[
 \sum_i(-1)^i\frac{i+\rho}{2}b_i
 =\frac{\rho}{2}\chi(E)+\frac12
 \underset{\substack{\rule{0.4pt}{4pt}\rule{6em}{0.4pt}\rule{0.4pt}{4pt}\\ |}}
 {\sum_{i=0}^{2n}(-1)^i i b_i}
\]
\[
 \sum_{i=0}^{n-1}\bigl[(-1)^i i b_i+
 (\deleted{\unsure{2|-2}}-1)^{2n-i}(2n-i)b_{2n-i}\bigr]+(-1)^n n b_n
\]
comme $b_i=b_{2n-i}$
\[
 2n\sum_{i=0}^{n-1}(-1)^i b_i+(-1)^n n b_n
 =\deleted{\unsure{2|+}}n\sum_{i=0}^{2n}(-1)^i b_i=n\chi(E)
\]
\[
 \sum_i(-1)^i\frac{i+\rho}{2}b_i
 =\deleted{n+\rho}\frac{n+\rho}{2}\chi(E) \qquad \text{donc}
\]
\[
 \delta(E)=\deleted{\unsure{\delta|\varepsilon}}\unsure{\delta|\varepsilon}(E)\,
 p^{\frac{n+\rho}{2}\chi(E)}
 \qquad \varepsilon(E)=\pm1
\]

Par suite l'équation fonctionnelle devient dans ce cas
\begingroup\csname tagsleft@true\endcsname\[
 L_{DE}(t)=\varepsilon(E)\left(-p^{\frac{n+\rho}{2}}\frac{1}{t}\right)^{\chi(E)}
 L_E(t^{-1}) \tag*{(2)}
\]\endgroup

Ou encore, si on admet que
\[
 \check E\deleted{\unsure{\otimes E|\simeq E}}\simeq E(\rho)
 \qquad \text{globalement sur }X,\ \text{\deleted{\unsure{d'où|donc}}}
\]
[il suffit que ce soit vrai sur le corps des \unsure{fonctions|fractions} de $X$, donc c'est \unsure{sans doute} tjrs vrai]
\[
 \deleted{\unsure{L_E|L_{\check E}}}D(E)\simeq
 \deleted{\check E\otimes\unsure{\mathbb Q_\ell(n)|\mathbb Q_\ell}}
 \check E(n)\underset{\substack{\diagup\\\text{shift}}}{[2n]}
 \simeq E(\rho+n)\underset{\substack{\diagup\\\text{shift.}}}{[2n]}
\]
et par suite
\[
 \deleted{L_{DE}\simeq L_{E(\rho+n)}=
 L_{\unsure{\mathbb R|D|R}f_* E(\rho+n)}}
\]
\[
 \mathbb R f_{\unsure{{}|*}}(D(E))\simeq
 \mathbb R f_*(E)(\rho+\deleted{\unsure{2|3}}n)[2n]\qquad \text{d'où}
\]
\[
 L_{DE}\unsure{(t)|{}}=L_{\mathbb R f_*D(E)}\unsure{(t)|{}}
 \simeq L_{\mathbb R f_*(E)}(p^{-(\rho+n)}t)
\]
\deleted{et il \inserted{\unsure{s'en|s'e}} \unsure{suit|vient}}

